\chapter{Limitaciones de los modelos y perspectivas}


\subsection{Limitaciones de los modelos}

Las membrana biológica de cualquier célula es un sistema altamente complejo al presentar diversos componentes como proteínas, lípidos y carbohidratos principalmente. Y la diversidad de cada uno de los componentes es amplia, es por esto que aislarlas en su contexto nativo no una tarea simple. Por esta razón surgen modelos que simplificados que permiten estudiar ciertos componentes aislados, como los empleados en esta tesis, las capas monomoleculares con sus ventajas mencionadas previamente en la \textbf{sección} \textbf{\ref{monocapas}} y las bicapas lipídicas. Aquí también es necesario mencionar que estos sistemas no representan la realidad de la dinámica los procesos biológicos de la membrana. Algunas de sus limitaciones son: 

\begin{itemize}
\item{Variedad de especies lipídicas: Las membranas biológicas contienen una gran variedad de especies lipídicas, mientras que los modelos artificiales suelen utilizar un número más limitado de lípidos.}

\item{Asimetría de los lípidos: Es difícil imitar la asimetría característica de la distribución de los lípidos en las bicapas de las membranas biológicas.}

\item{Otros componentes: Los modelos de membrana lipídica no tienen en cuenta otros componentes importantes de las membranas biológicas, como las proteínas y los elementos del citoesqueleto, que regulan la difusión de lípidos y proteínas a través de la membrana.}
\end{itemize}

El modelo CG Martini implica un grado de aproximación que surge de la pérdida de detalles atómicos y es probable que los valores de las propiedades e interacciones lípido-péptido y proteína-proteína estimadas presenten una desviación respecto a los valores medidos experimentalmente. En particular, algunos cambios locales en la estructura secundaria podría inducirse tras la unión de lípidos reguladores. Las restricciones diédricas de la columna vertebral de proteína estándar empleadas dentro del modelo Martini \citep{Risselada2009, Marrink2007} y la falta de direccionalidad del enlace de hidrógeno que surge del sistema de CG significa que los posibles cambios conformacionales inducidos por los lípidos no se capturan en los cálculos actuales. El modelo de agua polarizable \citep{Yesylevskyy2010} mejora la precisión de los cálculos de energía libre para dividir las cadenas laterales de aminoácidos en ambientes hidrofóbicos. Sin embargo, sigue habiendo una comprensión incompleta de su influencia en las interacciones lípido-agua y proteína-agua.


Estas limitaciones hacen que los modelos de membrana lipídica artificial no puedan reproducir completamente la complejidad y el funcionamiento de las membranas biológicas. Sin embargo, siguen siendo herramientas valiosas para el estudio de ciertos aspectos de la estructura y la dinámica de las membranas celulares.


\subsection{Perspectivas}

Experimentalmente sería conveniente realizar una optimización a la metodología de expresión y repurificación por HPCL. Para la expresión se puede aumentar la cantidad en masa de proteína fusión realizando cultivos en un bioreactor. Cambiar el medio de cultivo por ejemplo, un LB fortificado con nutrientes. La repurificación por HPLC se podría escalar a una columna semi-preparativa, de esta forma cargar mayor cantidad de péptido previo a la elución evitando la sobresaturación de la columna como sucede al usar una columna analítica.

Contando con una cantidad considerable de los péptidos entre 5 a 10 mg son varios los análisis biofísicos que se pueden realizar. Realizar experimentos de monocapas en el \ac{bam} para analizar la partición de los péptidos a determinadas fase lipídicas. Con la espectroscopia de absorción por reflexión infrarroja 
espectroscopía infrarroja de reflexión-absorción por modulación de la polarización (PM-IRRAS) se puede conocer en detalle la orientación de los péptidos en las monocapas \citep{Rabe2014}. La espectroscopia de dicroísmo circular de radiación orientada es otra técnica que se puede emplear para saber la orientación y conformación de los péptidos en \acf{mlvs} \citep{Muruganandam2013}.

Al inicio del desarrollo de esta tesis estuvimos realizando experimentos de micro calorimetría de barrido diferencial de los péptidos puros y mezcla lípido-péptido en \acf{mlvs} , sin embargo, por la limitada disposición de péptido puro no se continuaron. La finalidad de estas mediciones apuntaban a responder la pregunta cuál es la preferencia de partición de los péptidos en la fase gel o líquido cristalino de una membrana. También se realizaron mediciones preliminares por FTIR de los péptidos en MLV's enfocados a conocer el efecto de los lípidos sobre estructura secundaria.

Realizar simulaciones multiescala grueso y átomo integrado para obtener los perfiles de energía libre de las interacciones lípido-péptido \citep{Corey2019, Domanski2017, Hedger2016}. En primera instancia se realizarían simulaciones CG que permiten obtener una adecuada exploración de la interacción lípido-proteína e identificar los sitios de contacto, posteriormente convertir a modelo atomístico y correr simulaciones de átomo integrado para cuantificar las interacciones. Además, en el análisis de contactos lípido-péptido considerar diferentes partículas del lípido, es decir, incluir las colas hidrocarbonadas que también podrían interaccionar con los péptidos.

En humanos la membrana plasmática de eritrocitos tiene la composición lipídica en peso que se muestra  abajo \citep{Hildebrand1975, Bruce2002}. Realizar simulaciones de DM en un sistema más realístico en composición de la membrana podría aportar información más acertada sobre el efecto de la membrana en la conformación y orientación de la integrina.

%%%Thomas E. Andreoli et al., Membrane Physiology (1987), Table 1, Chapeter 27
\begin{table}[H]
\centering
%\caption{}
%\label{tab:my-table}
\begin{tabular}{ll}
\rowcolor[HTML]{EFEFEF}
\hline
Lípido                & \%   \\ \hline
colesterol            & 26.0 \\
fosfatidilcolina      & 17.5 \\
esfingomielina        & 16.0 \\
fosfatidiletanolamina & 16.6 \\
fosfatidilserina      & 7.9  \\
fosfatidilinositol    & 1.2  \\
otros                 & 12.5 \\ \hline
\end{tabular}
\end{table}

En el presente estudio se trabajó solo con los segmentos TM-IC, dado que ya se dispone de la estructura completa integrina $\alpha$IIb$\beta$3 \citep{Adair2023}. A nivel experimental también se podría contactar a los autores si acceden a compartir el plásmido de la estructura y abordar estudios biofísicos experimentales y por simulaciones de DM para responder la pregunta del efecto del ambiente lipídico en la conformación de la integrina completa.